\documentclass[../../../include/open-logic-section]{subfiles}

\begin{document}

\olfileid{sth}{ordinals}{opps}
\olsection{અનુગામી અને સીમા ક્રમવાચક સંખ્યાઓ}

આગળનાં થોડાં અધ્યાયોમાં આપણે ક્રમવાચક સંખ્યાઓનો ઘણો ઉપયોગ
કરીશું. તેથી કેટલાક સરળ ખ્યાલો રજૂ કરી લઈએ.

\begin{defn}
કોઈ પણ ક્રમવાચક સંખ્યા $\alpha$નો \emph{અનુગામી}
$\ordsucc{\alpha}
=\alpha \cup \{\alpha\}$ છે. કોઈ ક્રમવાચક સંખ્યા~$\beta$ માટે
$\ordsucc{\beta} = \alpha$ હોય, તો $\alpha$ને
\emph{અનુગામી} ક્રમવાચક સંખ્યા કહીએ. $\alpha$ ખાલી ન હોય અને
$\alpha$ અનુગામી ક્રમવાચક સંખ્યા પણ ન હોય, તો અને માત્ર તો તેને
\emph{સીમા} ક્રમવાચક સંખ્યા કહીએ.
\end{defn}
\noindent
આગળનું પરિણામ બતાવે છે કે \emph{અનુગામી}નો આ યોગ્ય ખ્યાલ છે:

\begin{prop}
કોઈ પણ ક્રમવાચક સંખ્યા $\alpha$ માટે:
\begin{enumerate}
	\item $\alpha \in \ordsucc{\alpha}$;
	\item $\ordsucc{\alpha}$ ક્રમવાચક સંખ્યા છે;
	\item $\alpha \in \beta \in
	\ordsucc{\alpha}$ થાય એવી કોઈ ક્રમવાચક સંખ્યા $\beta$ નથી.
	\end{enumerate}
\end{prop}

\begin{proof}
સ્પષ્ટ છે કે $\alpha \in \alpha \cup \{\alpha\} = \ordsucc{\alpha}$.
તે જ રીતે $\ordsucc{\alpha}$ ક્રમવાચક સંખ્યાઓનો પરંપરિત ગણ છે,
તેથી \olref[basic]{corordtransitiveord} મુજબ ક્રમવાચક સંખ્યા છે.
વળી $\alpha \in \beta \in \ordsucc{\alpha}$ અશક્ય છે,
કારણ કે ત્યારે $\beta \in \alpha$ અથવા $\beta = \alpha$,
જે \olref[basic]{ordordered}નો વિરોધ કરે છે.
\end{proof}
\noindent
આથી અનંતાતીત અનુમાન વડે સાબિતીની એક બીજી રીત પણ મળે છે:

\begin{thm}[સરળ અનંતાતીત અનુમાન]\ollabel{simpletransrecursion}
ધારો કે $\phi(x)$ એવું સૂત્ર છે કે:
\begin{enumerate}
	\item $\phi(\emptyset)$; અને
	\item કોઈ પણ ક્રમવાચક સંખ્યા $\alpha$ માટે, જો $\phi(\alpha)$ હોય, તો
	$\phi(\ordsucc{\alpha})$; અને
	\item જો $\alpha$ સીમા ક્રમવાચક સંખ્યા હોય અને $(\forall \beta \in
	\alpha)\phi(\beta)$ હોય, તો $\phi(\alpha)$.
\end{enumerate}
તો $\forall \alpha \phi(\alpha)$.
\end{thm}

\begin{proof}
આપણે પ્રતિવિધાન સાબિત કરીએ. ધારો કે એવી કોઈ ક્રમવાચક સંખ્યા છે
જેના માટે $\lnot\phi$; આવી લઘુતમ ક્રમવાચક સંખ્યા $\gamma$ લો.
ત્યારે કાં તો $\gamma = \emptyset$, અથવા કોઈ એવી $\alpha$ માટે
$\gamma = \ordsucc{\alpha}$ છે જેના માટે $\phi(\alpha)$;
અથવા $\gamma$ સીમા ક્રમવાચક સંખ્યા છે અને $(\forall
\beta \in \gamma)\phi(\beta)$. દરેક કિસ્સો સંબંધિત ધારણાથી
વિરોધાભાસ આપે છે.
\end{proof}
\noindent
આગળ એક અંતિમ સંકેત ઉપયોગી થશે:

\begin{defn}\ollabel{defsupstrict}
જો $X$ ક્રમવાચક સંખ્યાઓનો ગણ હોય, તો $\supstrict(X) = \bigcup_{\alpha
\in X} \ordsucc{\alpha}$.
\end{defn}
\noindent
અહીં ``lsub''નો અર્થ ``ન્યૂનતમ કડક ઉચ્ચસીમા'' છે.\footnote{કેટલાંક
પુસ્તકો તેના માટે ``$\text{sup}(X)$'' લખે છે. પરંતુ બીજા પુસ્તકો
``$\text{sup}(X)$''નો અર્થ ન્યૂનતમ \emph{અકડક} ઉચ્ચસીમા,
એટલે કે માત્ર $\bigcup X$, એવો કરે છે. જો $X$માં મહત્તમ સભ્ય
$\alpha$ હોય, તો આ બે ખ્યાલ જુદા પડે: ન્યૂનતમ \emph{કડક}
ઉચ્ચસીમા $\ordsucc{\alpha}$ છે, જ્યારે ન્યૂનતમ \emph{અકડક}
ઉચ્ચસીમા માત્ર $\alpha$ છે.} આગળનું પરિણામ તેનો ખુલાસો કરે છે:

\begin{prop}
જો $X$ ક્રમવાચક સંખ્યાઓનો ગણ હોય, તો $\supstrict(X)$ એ $X$ની
દરેક ક્રમવાચક સંખ્યા કરતાં મોટી હોય એવી ન્યૂનતમ ક્રમવાચક સંખ્યા છે.
\end{prop}

\begin{proof}
$Y = \Setabs{\ordsucc{\alpha}}{\alpha \in X}$ લો, જેથી
$\supstrict(X) = \bigcup Y$. ક્રમવાચક સંખ્યાઓ પરંપરિત હોવાથી અને
ક્રમવાચક સંખ્યાનો દરેક સભ્ય પણ ક્રમવાચક હોવાથી $\supstrict(X)$
એ ક્રમવાચક સંખ્યાઓનો પરંપરિત ગણ છે; તેથી
\olref[basic]{corordtransitiveord} મુજબ તે ક્રમવાચક સંખ્યા છે.

જો $\alpha \in X$ હોય, તો $\ordsucc{\alpha} \in Y$, તેથી $\ordsucc{\alpha}
\subseteq \bigcup Y = \supstrict(X)$, અને તેથી $\alpha \in
\supstrict(X)$. આમ $\supstrict(X)$ એ $X$ની દરેક ક્રમવાચક
સંખ્યા કરતાં કડક રીતે મોટો છે.

ઊલટું, જો $\alpha \in \supstrict(X)$ હોય, તો કોઈ $\beta \in X$
માટે $\alpha \in
\ordsucc{\beta} \in Y$, તેથી $\alpha \leq
\beta \in X$. એટલે $\supstrict(X)$ એ $X$ની \emph{ન્યૂનતમ}
કડક ઉચ્ચસીમા છે. ખરેખર, આપેલા ગણની બધી ક્રમવાચક સંખ્યાઓથી
મોટી કોઈ પણ ક્રમવાચક સંખ્યા આ બધાં સભ્યોને સમાવે છે.
\end{proof}

\end{document}
